\documentclass{article}
\usepackage{amsmath}
\begin{document}
\section*{Transporting a risk warrant through a changed observation map}
\paragraph{Scope.}
This follow-up to PR 851 pairs the two pinned radio configurations on the
same upstream event coordinates. It separates three experiments:
a shared distribution on actual attempted sends, fresh separately asserted
outcome marginals, and two calibrated experiments with a certified bound on
how much their upstream distributions differ. No operational bound on that
difference and no held-out evidence are supplied. The bounded comparison
assumes only the two sets of exact marginals and a bound between the two
otherwise unknown laws; it does not take the full old distribution as known.

\paragraph{Source and preserved query.}
Let $U$ indicate a two-byte attempted send, and $V$ a one-byte attempted
settlement of the old keyed report. The cap-2 and cap-3 outcome maps on
$(U,V)$ are $(U(1-V),V)$ and $(U,V)$. Paired native rows show the same
incremental old-loss event $\{00\}$ for the admitted branch under each
map, and no old-duty loss in the withheld branch. Unrelated service
survives in all pairs. With row-0 admission probability $a$ and independent
randomization, the same upstream law thus preserves old risk exactly,
regardless of whether it puts mass on the new safe output $(1,1)$.
The optional completion event is different: for the same law and $a$,
cap 3 loses exactly $a\Pr(U=V=1)$ optional completions, because the newly
delivered two-byte job leaves no room for the optional packet.
Equating old-risk queries does not give service equivalence.

If old outcome marginals are $p=\Pr(X=1)$ and $q=\Pr(Y=1)$, the old
attempted-send distribution must be
\[
   (1-p-q,\ q-s,\ p,\ s),\qquad 0\le s\le q.
\]
In particular old risk is $a(1-p-q)$. For a shared law the new marginals
are $P=p+s,\ Q=q$. Measuring both $P=p$ and $Q=q$ after the upgrade is
consistent with a shared law only if $s=0$; the new pair may be reachable
without positive probability. Unknown dependence on the old experiment's
unresolved split $s$ does not affect either experiment's old risk at zero
drift, but still affects optional service if $P$ is not measured.

\paragraph{General exact recalibration.}
Write the new outcome marginals as $P,Q$ and the new attempted-send law as
$(1-P-Q+t,\ Q-t,\ P-t,\ t)$ with
$\max(0,P+Q-1)\le t\le\min(P,Q)$. Set $d=P-p$ and $e=Q-q$.
Fix $t$. Over all old laws consistent with $p,q$, the least total
variation distance to the new law is
\[
  D(t)=\max\bigl(|t-d-e|,\ |d-t|,\ |e|\bigr).
\]
Indeed the first and third signed cell changes are $t-d-e$ and $d-t$.
The changes in the two remaining cells sum to $e$, so their absolute
values sum to at least $|e|$. Choose $s=\min(q,t)$ to attain this lower
bound: if $t\le q$, the fourth cell does not change; if $t>q$, both
changes in the $V=1$ block are nonnegative. The three signed quantities
sum to zero, so half their absolute sum equals their largest absolute
value. This is also the coarse projection onto the three atoms
$\{00\},\{10\},\{01,11\}$. The projection loses no distance information
for this \emph{existential} comparison with freely chosen old split.
A separately known old split can yield a stronger, different test.

Assume $p,q,P,Q\in[0,1]$ and $\delta\ge0$.
The four marginal observations and an upstream drift warrant $\delta$
are jointly consistent exactly when $p+q\le1$, $|e|\le\delta$, and
\[
 \max(0,P+Q-1,d-\delta,d+e-\delta)
 \le
 \min(P,Q,d+\delta,d+e+\delta).
\]
These are the endpoints of the compatible $t$ interval.
Let $t_{\max}$ be its right endpoint. The sharp new worst-case risk
is $a(1-P-Q+t_{\max})$; optional completions have expectation
$a(1-P)$. An empty interval is a failed joint calibration, not a zero
risk verdict. At fixed half marginals, risk is exactly
$a\min(\delta,1/2)$ in the worst case. At zero drift and consistent
marginals, the old risk answer carries over unchanged; no new
transport solve is needed. Re-solving with the same upstream model
gives the same answer. The simplification does not carry to new
loss preimages, departure feedback changing the traffic law outside
the declared bound, or uncertainty in how the attempted workload was
measured.

\paragraph{Checks.}
The paired comparison reads fresh native source results and checks both
admission arms on all four attempt patterns under both caps.
The independent rational sweep enumerates all $969^2$ pairs of
sixteenth-unit laws, varying both sets of marginals and each drift
budget from zero to one. Lean checks the lower bound and attainment in
the unobserved split and the exact absolute-value interval.
\end{document}
